\scanpage{23}
\mspage{20}

F6) Soient $I,J\in S$, \inserted{$x\in L$,} tels que $x\in\delta I$, $x\in\delta J$, et que $I,J$ définissent la $=$ branche i.e. $\exists y\in I^\circ\cap J^\circ$, \inserted{$y\ne x$,} tel que $I_{x,y}=J_{x,y}$. Alors $\exists K\in S$ tel que
\[
x\in\delta K,\ I^\circ\cup J^\circ\subset K^\circ.
\]

En d'autres termes, l'ens des \inserted{$I^\circ$,} $I\in S$ tels que $x\in\delta I$, et que $I$ appartienne à une branche donnée \inserted{en $x$,} est filtrant croissant pour l'inclusion.

\margin{\uline{NB} On n'exige pas que les autres extrémités \inserted{$\deleted{\unsure{\delta}}I,\deleted{\unsure{\delta}}J$} \unsure{de $I$} soient dans $K$. Mais \deleted{\unsure{il est clair}} que celles de ces extrémités qui sont régulières \unsure{soient} \deleted{\unsure{dans}} $K$ \unsure{Car conséquence de la propriété de dilatation|Car conséquence de la propriété d'extension} (Cor. 2 p. 7).}

Ou encore (en connexion avec F3)

Soient $I,J\in S$, $x\in L$ tels que $x\in\delta I\cap\delta J$. Alors

a) Pour qu'il existe $y\in I^\circ\cap J^\circ$ tel que
\[
I_{x,y}=J_{x,y},
\]
il f. et s. qu'il existe $K\in S$ tel que $x\in\delta K$, $I\cup J\subset K$

b) Si cette condition \unsure{n'est} vérifiée alors $\exists u\in I^\circ$, $v\in J^\circ$ tels que
\[
I_{x,u}\cap\unsure{I_{x,v}|J_{x,v}}=\{x\}.
\]

\mspage{22}

\rotatebox{90}{\reflectbox{\unsure{EXTRA STRONG}}}
\rotatebox{90}{\reflectbox{\unsure{EXTRA STRONG}}}
